% Part: sets-functions-relations
% Chapter: arithmetization
% Section: reflections

\documentclass[../../../include/open-logic-section]{subfiles}

\begin{document}

\olfileid{sfr}{arith}{ref}
\olsection{ځينې فلسفي غورونه}

فني خبرې تر دې ځايه بس دي۔ خو هغو څه ترلاسه کړل؟

ښه، دا خبره تقريباً له اختلافه پاکه ده چې هغو موږ ته څه {ښکلې}
سوچه رياضي راکړه۔ برسېره پر دې، ځينې ژورې تصوري لاسته راوړنې هم
وې۔ دا پېژندنه ژوره وه چې د بشپړتيا خاصيت د حقيقي او ناطقو عددونو
تر منځ بنسټيز توپير بيانوي۔ همدارنګه د ډېډېکېنډ د پرېکونونو په توګه
د حقيقي عددونو ښکاره جوړښت د تحليل موضوع ته کلک بنسټ ورکوي۔ پوهېږو
چې د \emph{بشپړ مرتب ميدان} تصور بې تناقضه دے، ځکه پرېکونونه کټ
مټ همداسې يو ميدان جوړوي۔

له دې ټولو سره، د دغو لاسته راوړنو په اړه بايد څو تحفظات څرګند کړو۔

لومړے، دا څرګنده نۀ ده چې د پرېکونونو په توګه د حقيقي عددونو ګڼل
د هغو د پېژندل شوو (ښايي نامتناهي) اعشاري پراختياوو په توګه تر ګڼلو
څومره \emph{زيات} دقيق دي۔ د حقيقي عددونو دا وروستے «جوړښت» د
کوشي تسلسلونو له لارې د حقيقي عددونو جوړښت ته څه ورته والے لري؛ خو
په حقيقت کښې رياضي پوهان د اوولسمې پېړۍ له پيله په بنسټيز ډول پرې
پوهېدل (\olref[cauchy]{sec} وګورئ)۔ په دقت کښې رښتينے زياتوالے له
دې پېژندنې راغے چې حقيقي عددونه د بشپړتيا خاصيت لري؛ يوازې دا وړتيا
چې حقيقي عددونه د ځانګړو سټونو په توګه جوړ کړو، ښايي پخپله دومره
په زړه پورې نۀ وي۔

دا لا هم لږ څرګنده ده چې د صحيح يا ناطقو عددونو (ډېرې اسانه) سټي
جوړونه په دغو ډګرونو کښې دقت زياتوي۔ دلته يوه ساده خبره د يادولو وړ
ده۔ صحيح عددونه مو د طبيعي عددونو د مرتبو جوړو د هم ارزښتۍ ټولګيو
په توګه \emph{جوړ} کړل، بيا مو ناطق عددونه د صحيح عددونو د مرتبو
جوړو د هم ارزښتۍ ټولګيو په توګه جوړ کړل، او بيا مو حقيقي عددونه د
ناطقو عددونو د سټونو په توګه جوړ کړل؛ خو سمدستي دا جوړښتونه
\emph{هېروو}۔ په ځانګړي ډول هېڅوک به کله هم ونۀ غواړي چې په يوه
رياضياتي ثبوت کښې دا جوړښتونه \emph{وکاروي} (البته د دې ثبوت پرته
چې جوړښتونه لکه څنګه چې ټاکل شوي وو، هماغسې چلېږي)۔ د يوه حقيقي
عدد په اړه نېغ په نېغه خبرې کول د طبيعي عددونو د سټونو د سټونو د
سټونو د سټونو د سټونو د سټونو د کوم سټ په اړه تر خبرو ډېر اسانه دي۔
%(ځکه حقيقي عددونه د ناطقو عددونو د سټونو په توګه جوړ شول؛ او دا د صحيح عددونو د مرتبو جوړو د هم ارزښتۍ ټولګيو، يعنې د صحيح عددونو د سټونو د سټونو د سټونو په توګه جوړ شول؛ او داسې نور۔)
%2/3 = [2,3]
%	يعنې {<2,3>, ... }
%	{ {{2},{2,3}}, ... }
%
%حقيقي عددونه
%	د ناطقو عددونو سټونه دي
%
%ناطق عددونه
%	د صحيح عددونو د سټونو د سټونو سټونه دي
%
%صحيح عددونه
%	د طبيعي عددونو د سټونو د سټونو سټونه دي
	
تر ټولو زيات شک په دې دے چې دا تعريفونه موږ ته وايي چې صحيح، ناطق
يا حقيقي عددونه په \emph{ماوراءالطبيعتي معنا څه دي}۔ يعنې په دې شک
دے چې حقيقي عددونه (مثلاً) ځينې ځانګړي سټونه (د سټونو د سټونو
\ldots) \emph{دي}۔ د دې نظر پر وړاندې ستر خنډ دا دے چې جوړښت په
ډېرو بېلابېلو لارو کېدے شو۔ د حقيقي عددونو په حالت کښې ځينې په
رښتيا په زړه پورې بېلابېل جوړښتونه شته (\olref[cauchy]{sec}
وګورئ)۔ خو د بېلابېلو جوړښتونو د ترلاسه کولو يوه ډېره ساده لار دا
ده: لکه په \olref[sfr][rel][ref]{sec} کښې، مرتبې جوړې مو لږ په
بله توګه تعريفولے شوې؛ که مو د مرتبې جوړې دا بديل تصور کارولے وے،
زموږ جوړښتونو به کټ مټ همدومره سم کار کړے وے، خو په پاے کښې به مو
بېل څيزونه ترلاسه کړي وو۔ له دې امله د صحيح، ناطقو او حقيقي عددونو
ډېر سيال سټي جوړښتونه شته۔ اوس دا ادعا چې صحيح عددونه (مثلاً)
\emph{دا} سټونه دي، نۀ \emph{هغه}، يوازې خپلسرې (او شرموونکې) به
وي۔ (لکه په \olref[sfr][rel][ref]{sec} کښې، دا د هغه استدلال يوه
بېلګه ده چې \citealt{Benacerraf1965} مشهور کړ۔)

يو بل ټکی هم د راپورته کولو وړ دے: زموږ په جوړښتونو کښې يو ډېر
\emph{عجيب} څيز شته۔ موږ په طبيعي عددونو پيل وکړ۔ بيا صحيح عددونه
جوړوو، او «د صحيح عددونو $0$»، يعنې
$ \equivrep{0,0}{\Intequiv}$، جوړوو۔ خو
$0 \neq \equivrep{0,0}{\Intequiv}$۔ په حقيقت کښې زموږ د جوړښتونو
له مخې \emph{هېڅ} طبيعي عدد صحيح عدد نۀ دے۔ خو دا د شهودي پوهې
سخت خلاف ښکاري۔ همدارنګه په \olref[sfr][set][imp]{sec} کښې مو له
ډېر استدلال پرته ادعا وکړه چې $\Nat \subseteq \Rat$۔ که جوړښتونه
موږ ته په دقيق ډول وايي چې عددونه \emph{څه} دي، نو دا ادعا په
څرګند ډول ناسمه وه۔

نو چې يو ګام شاته ودرېږو، کوم ځاے ته رسېږو؟ په ساده سټ نظريه کښې
په کار کولو او طبيعي عددونو ته په ورکړل شوو څيزونو کتلو، صحيح، ناطق
او حقيقي عددونه د ځينو سټونو په توګه \emph{کارولے} شو۔ په دې معنا
د دغو څيزونو نظريې په سټ نظريه کښې \emph{ورګډولے} شو۔ خو د دې
ورګډولو فلسفي ارزښت دومره ساده نۀ دے۔

البته دا هېڅ د بحث وروستۍ خبره نۀ ده!{} ټکی يوازې دا دے: د دې د
ښودلو دپاره چې د حقيقي عددونو سټي جوړونه ژور فلسفي ارزښت \emph{لري}،
يو څه نور \emph{فلسفي} استدلال ته اړتيا ده۔

%برسېره پر دې، د دې ټول څپرکي په اوږدو کښې مو طبيعي عددونه يوازې
%د راکړل شوو څيزونو په توګه ګڼلي دي۔ خو له هغو سره څه کولے شو؟
%دا د بل څپرکي موضوع ده۔

\end{document}
